\section{Lefebvre before Écône}

Marcel Lefebvre's authority among the priests who later gathered around him rested on a long ecclesiastical career, not merely on opposition to a new missal. Born at Tourcoing in 1905, ordained a priest in 1929, and professed in the Congregation of the Holy Ghost, he served as a missionary and seminary rector in Gabon before transfer to Senegal. He became apostolic vicar of Dakar and was consecrated bishop in 1947. Pius XII appointed him apostolic delegate for French-speaking Africa in 1948 and archbishop of Dakar in 1955. He briefly governed the Diocese of Tulle in 1962 before election as superior general of the Holy Ghost Fathers. He resigned that office in 1968 amid conflict over the congregation's postconciliar direction.

This career furnished three resources. Mission government gave him administrative experience and a network of clergy. Episcopal office gave him the sacramental capacity that would later make independent priestly and episcopal formation possible. His years in Rome and Africa also made his eventual opposition appear to supporters as the resistance of an experienced churchman rather than the protest of a marginal local organizer. The same biography sharpened the conflict for the Holy See: a former apostolic delegate and religious superior was not plausibly unaware of the meaning of papal warnings or canonical mission.

\subsection{Council participant and organized minority}

At the Second Vatican Council Lefebvre took part in the conservative episcopal network usually called the \emph{Coetus Internationalis Patrum}. Its members tried to amend conciliar texts and resisted formulations they believed weakened inherited teaching, especially on religious liberty, ecumenism, collegiality, and the Church's relation to the modern state. This history should not be simplified into a claim that Lefebvre rejected every conciliar document at the moment of promulgation. The documentary question of which acts he signed is distinct from his later interpretation and rejection of particular teachings. What became institutionally decisive was his conviction that the postconciliar settlement represented more than imprudent policy: he increasingly described it as a doctrinal and spiritual crisis.

Liturgical reform concentrated that diagnosis. The revised rites were promulgated by competent papal authority, but Lefebvre judged their theology, ecumenical orientation, manner of celebration, and practical fruits together. The resulting controversy was never only about whether Latin, chant, eastward celebration, or an older calendar should survive. For him the older liturgy became a visible rule of belief and an instrument for forming priests against a wider ecclesial transformation. For Roman authorities, however, fidelity to tradition could not be set against obedience to the living magisterium and the lawfully promulgated liturgy.

\begin{studybox}{Three questions that later became one conflict}
\textbf{Formation:} Could a seminary form priests according to older manuals, disciplines, and liturgical books while rejecting major features of the postconciliar settlement?\par
\textbf{Doctrine:} Did disputed conciliar teachings develop Catholic doctrine or contradict it?\par
\textbf{Authority:} Who could judge that question in a way binding on clergy and faithful, and could an emergency justify ministry against ordinary and papal commands?
\end{studybox}

The eventual SSPX answer joined the questions: sound priestly formation required resistance, and resistance required institutions. The Roman answer also joined them: a formation project could not claim Catholic mission while treating its own judgment of necessity as superior to hierarchical governance. This structure was already present before the penal confrontations of 1976 and 1988.
